\documentclass{article}

% Do not modify neurips_2026.sty.
\PassOptionsToPackage{numbers, compress}{natbib}
\usepackage[preprint]{neurips_2026}

% Counteracts the style's flushbottom, which stretches short pages into visible gaps.
\raggedbottom

\usepackage[utf8]{inputenc}
\usepackage[T1]{fontenc}
\usepackage{hyperref}
\usepackage{url}
\usepackage{booktabs}
\usepackage{amsfonts}
\usepackage{nicefrac}
\usepackage{microtype}
\usepackage{xcolor}

\title{TradeCopilot: A Personalized AI Trading Assistant with Structured User Memory}

\begin{document}

\maketitle

\section{Project objective}

We propose TradeCopilot, a personalized AI trading assistant for self-directed investors in US
equities. It researches a stock or a portfolio on demand, checks the main claims of its own
analysis against primary sources where possible, uses structured user memory to personalize its
recommendations, and supports user-controlled execution modes, including approval-gated and
bounded autonomous operation.

Our objective is to evaluate whether structured user memory improves profile fit, constraint
adherence, and decision consistency relative to a memoryless agent and to weaker profile-memory
baselines on scenarios where memory should matter.

\section{Motivation and challenges}

Self-directed investors often lose money not from a lack of information but from breaking their
own rules: they oversize positions, chase news, and drift past the risk limits they set for
themselves. At the same time, thorough research across filings, news, and market data
takes hours that most people do not spend. General chat assistants can answer a question about a
stock, but they do not know the user, they cannot be held to the user's own rules, and they cannot
be trusted anywhere near an actual trade. As more people turn to AI for financial questions, the
stakes rise: models state wrong numbers with full confidence, and no one checks them.

Consider Daniel, who uses an AI chat assistant to research stocks after work. He recently
lowered his single-stock limit from 10\% to 5\% of his portfolio and rejected a short-term
trade because he could not monitor the market during work. During his lunch break, a company
he owns reports strong earnings, and he asks: ``Should I buy more?'' He wants an answer that
reflects what he has already explained, but still finds himself repeating his preferences
and checking for outdated assumptions. TradeCopilot is designed to carry this context across
conversations. It retains his latest confirmed limits, records why he accepted or rejected
earlier proposals, and uses that history alongside current holdings and fresh evidence.
When Daniel asks about buying more, it would check how the purchase fits his updated limit
and explain whether the opportunity suits his investment approach. He could spend less time
repeating himself and correcting unsuitable advice, and more time considering his options
with an assistant that builds on its understanding of him.

The first challenge is memory that remains available at decision time. Existing agent harnesses and LLM trading agents already have memory
\citep{memgpt,finmem,tradinggpt}, so building another assistant that remembers the user is not
the main issue. The harder problem is representing a user's profile in a form the system can
reuse consistently across decisions. Risk limits and execution permissions should remain available at
action time instead of being left only in prompt text, while softer preferences should still shape the
assistant's judgment. Our evaluation addresses this by comparing memory levels under the same
execution safeguards and measuring profile fit and rule adherence.

The second challenge is verification and evaluation in a domain without ground truth. Claims about
a company are volatile and contradictory across sources, and language models invent plausible
figures, so every claim that shapes a decision has to be checked against primary sources before it
is used.
Measuring whether that verification works requires ground truth that we have to author ourselves,
which is why our evaluation includes deliberately planted false claims. The problem extends to the
overall evaluation: because an LLM's training data contaminates any replay of historical periods,
backtesting is not a valid measurement for our system, so success has to be assessed through frozen
snapshots and rubric-based judgment, with agreement between human and automated scoring reported.

The third challenge is safety and trust. Unrestricted autonomous trading by a language model is
not acceptable~\citep{amodei}, so execution must remain bounded by user-selected permissions and
code-level constraint checks rather than by the prompt alone.

\section{Potential solutions}

\subsection{System overview and architecture}

We are building a co-pilot for self-directed investors in US equities. It researches on demand
and helps the user stick to their own trading rules. The user asks about a stock or a portfolio,
the agent gathers evidence across market data, filings, news, social sentiment, and macro
indicators, and the system then either requests approval or acts within a bounded autonomous mode
defined by previously authorized settings. The intended contribution is to keep user-specific constraints and preferences
available in structured form, so that decisions depend less on ad hoc prompt context and more on
persistent state. Relative to existing tools, the proposed system combines on-demand analysis,
source checking, user memory, and trading access in a single workflow. Bounded autonomous
execution is supported as an execution mode, but the evaluation in this proposal focuses on memory
rather than on execution-mode comparison.

The system has three tiers: a user interface for interaction and portfolio views, a backend that
manages sessions and memory, and an agent layer that accesses external data and execution tools.
The agent draws on market data, filings, news, social signals, and macroeconomic information.
Each observation is stored with source and time metadata so that stale or weakly supported
evidence can be flagged before it is used.

\subsection{Structured user memory}

A main contribution of the project is the memory subsystem. Unlike typical agent memory, which
only informs what the model says, our memory is organized around how user state is stored and
reused. It has three layers. Structured profile memory stores durable user state as typed fields
rather than free text: risk band, position caps, sector exclusions, horizon, watchlist, source
preferences, and execution permissions. Episodic decision memory records each proposal, the
evidence behind it, the user's action and reason, and what actually happened afterwards.
Working memory holds temporary run-specific state such as intermediate findings and claim-check
results. Conversational preferences can be promoted into the profile, but high-impact changes are
confirmed by the user before they become hard constraints. This structure differs from purely
conversational memory in three ways. First, important user state is represented in typed fields
that can be queried directly. Second, the journal is grounded in outcomes such as fills,
rejections, and profit or loss rather than in conversation alone~\citep{reflexion}. Third, the
profile and journal can be included or withheld separately during evaluation, which makes it
possible to test how each memory component affects system behavior.

\subsection{Decision and execution flow}

The agent plans, acts, and observes in a loop, and two checks are mandatory before a proposal
reaches the user. The first is verification: the main claims in the draft recommendation are
extracted and sorted into verifiable, opinion, or unverifiable. Verifiable claims are checked
against filings first, including the reported numbers, and then against news sources~\citep{factscore}. Each verifiable claim
is marked supported or contradicted, and contradicted claims are either removed or clearly
flagged. The second is a constraint check against the user's profile. After both pass, the system
either requests approval or executes under a bounded autonomous mode with previously authorized settings.

The agent decides the position size itself, using the user's profile as context. Proposing an
order is unrestricted, but execution remains subject to user authorization and hard typed
constraints such as position size, sector exclusions, and execution permissions. Softer
preferences continue to guide the recommendation without acting as blocking rules. Durable memory
remains user-correctable, and high-impact inferred changes require confirmation.

\section{Experimental dataset description and evaluation plan}

\subsection{Experimental Dataset}

We will build a dataset of 10 scenarios using frozen evidence from market, filing, news, social,
and macro sources. Each scenario will specify the instrument,
date, user request, initial portfolio, user profile, relevant history, and acceptance criteria.
Evidence will be collected once and reused across conditions, with its sources and timestamps
preserved. Historical context will contain only information available before the task.

Six scenarios will represent ordinary use, with two cases in each of three categories: earnings,
news, and portfolio adjustments. Some will involve updated preferences
or previous decisions, allowing us to test whether memory is used correctly. The remaining four
will be adversarial: two with deliberately false claims, one encouraging a violation of the
user's constraints, and one involving a tool failure.

We will check the reference facts and scoring criteria before evaluation, using separate
development examples to refine them. Each scenario will define acceptable outcomes rather than
one exact answer. Asking for clarification or reporting insufficient evidence may be appropriate,
but will be recorded separately from completing a recommendation.

\subsection{Evaluation Plan}

The evaluation will examine whether memory improves system behavior under fixed execution
safeguards, and what these benefits cost. Verification is assessed within the controlled scenario
set and is intended to measure claim-checking performance in that setting rather than to
establish a general guarantee of factual correctness.

For evaluation, we will vary the memory condition available to the same agent:

\begin{itemize}
  \item \textbf{No memory:} no persistent user memory.
  \item \textbf{Prose profile memory:} the user profile provided as ordinary text, representing a
        weak-memory baseline.
  \item \textbf{Structured profile memory:} the same profile information provided as typed
        constraints and preferences, without the decision journal.
  \item \textbf{Full memory:} structured profile memory together with the relevant decision journal.
\end{itemize}

All conditions will receive the same current request, market evidence, and portfolio state.
We will keep the model, tools, other instructions, and resource limits fixed, and reset the
scenario before each run. Execution safeguards will remain identical across memory levels.
Initial proposals will be scored before constraint-check feedback, separating the agent's use
of memory from the system's enforcement of rules.

The full system will run on all 10 scenarios. The memory comparison is a focused ablation on a
subset selected in advance to cover constraint use, preference changes, and previous decisions.
We will repeat those subset comparisons to assess consistency while keeping the pilot manageable.

We will report four groups of measures:

\begin{itemize}
  \item \textbf{Task success and constraint adherence:} whether the scenario's acceptance
      criteria are met, how often initial proposals violate hard user constraints, and whether
        the execution check blocks invalid orders without rejecting valid ones.
  \item \textbf{Verification quality:} how many planted false claims are detected and how
        often true claims are incorrectly flagged.
  \begin{samepage}
  \item \textbf{Recommendation quality:} a predefined rubric covering evidence coverage,
        reasoning, risk disclosure, profile fit, and actionability. A judge from a different
        model family~\citep{zheng} will receive the reference evidence and user context,
        without seeing the memory condition. Human reviewers will check 20\% of outputs
        across conditions, plus disputed cases, and agreement will be reported.
  \par
  \end{samepage}
      \item \textbf{Operational performance:} tokens, cost, and latency.
\end{itemize}

Programmatic checks will assess measurable requirements wherever possible. Trace logs will
support failure analysis and help distinguish missing evidence, incorrect memory use, and reasoning
errors. Confirmed failures will become regression cases for later versions.

Results will include per-scenario differences and variation across repeated runs. This small
study evaluates assistant behavior rather than investment returns; frozen evidence improves
comparability but does not eliminate possible knowledge from model training.

\nocite{liagents}
\bibliographystyle{plainnat}
\bibliography{refs}

\end{document}
